\documentclass[12pt,a4paper]{article}

% 中文支持
\usepackage{xeCJK}
\setCJKmainfont{Songti SC}
\setCJKsansfont{Heiti SC}
\setCJKmonofont{Heiti SC}

% 数学包
\usepackage{amsmath,amssymb,amsthm}
\usepackage{bm}
\usepackage{mathtools}
\usepackage{booktabs}

% 页面设置
\usepackage{geometry}
\geometry{left=2.5cm,right=2.5cm,top=2.5cm,bottom=2.5cm}

% 定理环境
\newtheorem{theorem}{定理}
\newtheorem{lemma}{引理}
\newtheorem{proposition}{命题}
\newtheorem{assumption}{假设}
\newtheorem{remark}{注}

% 自定义命令
\newcommand{\vect}[1]{\mathbf{#1}}
\newcommand{\mat}[1]{\mathbf{#1}}
\newcommand{\norm}[1]{\|#1\|_2}
\newcommand{\tr}{\text{tr}}
\newcommand{\Real}{\text{Re}}
\DeclareMathOperator{\rank}{rank}
\DeclareMathOperator*{\argmin}{arg\,min}
\DeclareMathOperator*{\argmax}{arg\,max}

\title{\textbf{Cell-Free ISAC 系统鲁棒波束成形优化\\数学推导}}
\author{}
\date{}

\begin{document}

\maketitle

%==============================================================================
\section{原始问题}
%==============================================================================

原始 (P1) 问题的决策变量为 per-AP 波束 $\vect{w}_{m,k} \in \mathbb{C}^{N_t}$、感知协方差 $\mat{Z}_m \in \mathbb{H}_+^{N_t}$、AP 选择指示 $b_{mp} \in \{0,1\}$：
\begin{align}
\min_{\substack{\{\vect{w}_{m,k}\}_{m \in \mathcal{M}, k \in \mathcal{K}}, \{\mat{Z}_m\}_{m \in \mathcal{M}}, \\ \{b_{mp}\}_{m \in \mathcal{M}, p \in \mathcal{P}}}} \quad 
& \sum_{m=1}^{M} \Big( \sum_{k=1}^{K} \|\vect{w}_{m,k}\|^2 + \tr(\mat{Z}_m) \Big) \tag{P1-objective} \\
\text{s.t.} \quad 
& \text{SINR}_{S,p} := \frac{|\vect{g}_p^H \mat{Z} \vect{g}_p|}{\sigma_s^2} \geq \gamma_S^{\text{PoD}}, \quad \forall p \in \mathcal{P} \tag{P1-C2} \\
& \tr(\mat{J}_p^{\text{data}}) \geq \Gamma_{\text{Track}, p}, \quad \forall p \in \mathcal{P} \tag{P1-C3} \\
& \sum_{m=1}^{M} b_{mp} = N_{\text{req}}, \quad \forall p \in \mathcal{P}^{\text{active}} \tag{P1-C4} \\
& \sum_{k=1}^{K} \|\vect{w}_{m,k}\|^2 + \tr(\mat{Z}_m) \leq P_{\max}, \quad \forall m \in \mathcal{M} \tag{P1-C5} \\
& \mat{Z}_m \succeq \mat{0}, \quad \forall m \in \mathcal{M} \tag{P1-C6} \\
& b_{mp} \in \{0, 1\}, \quad \forall m \in \mathcal{M}, p \in \mathcal{P} \tag{P1-C7}
\end{align}

通信 SINR (P1-C1) 与最坏情况 SINR 定义为（**注意：分式结构与半无限不确定性是核心非凸源**）：
\begin{align}
\tag{P1-C1} \text{SINR}_k^{\text{wc}} &:= \min_{\|\Delta\vect{h}_k\| \leq \epsilon_h \|\hat{\vect{h}}_k\|} \frac{|(\hat{\vect{h}}_k + \Delta\vect{h}_k)^H \vect{w}_k|^2}{\sum_{j \neq k} |(\hat{\vect{h}}_k + \Delta\vect{h}_k)^H \vect{w}_j|^2 + \sigma_c^2} \geq \gamma_k, \quad \forall k
\end{align}

其中 $\vect{w}_k \triangleq [\vect{w}_{1,k}^T, \vect{w}_{2,k}^T, \ldots, \vect{w}_{M,k}^T]^T \in \mathbb{C}^{MN_t}$ 为所有 AP 到用户 $k$ 的协作合并波束向量（**注意**：决策变量是 per-AP 的 $\vect{w}_{m,k}$，$\vect{w}_k$ 由其堆叠得到）。

%==============================================================================
\section{协方差提升形式 (P2) — 等价 lifted 形式}
%==============================================================================

为后续半定松弛做准备，将 (P1) 中向量变量提升为协方差矩阵，得到 \textbf{(P2)}：
\begin{align}
\min_{\substack{\{\mat{W}_k\}_{k\in\mathcal{K}}, \mat{Z} \in \mathbb{H}_+^{MN_t}, \\ \{b_{mp}\}_{m\in\mathcal{M},p\in\mathcal{P}}}} \quad
& \sum_{k=1}^K \tr(\mat{W}_k) + \tr(\mat{Z}) \tag{P2} \\
\text{s.t.} \quad
& \text{(P1-C1) 改写为 (P2-C1)} \tag{P2-C1} \\
& \tr(\vect{g}_p \vect{g}_p^H \mat{Z}) \geq \gamma_S^{\text{PoD}} \sigma_s^2, \quad \forall p \tag{P2-C2} \\
& \tr(\mat{F}_p \mat{R}_X) \geq \Gamma_{\text{Track},p}, \quad \forall p \tag{P2-C3} \\
& \tr(\mat{E}_m \mat{R}_X) \leq P_{\max}, \quad \forall m \tag{P2-C4} \\
& \mat{W}_k \succeq \mat{0}, \quad \forall k \tag{P2-C5} \\
& \mat{Z} \succeq \mat{0} \tag{P2-C6} \\
& \rank(\mat{W}_k) = 1, \quad \forall k \tag{P2-C7} \\
& b_{mp} \in \{0,1\}, \quad \forall m, p \tag{P2-C8}
\end{align}

其中：
\begin{itemize}
    \item $\mat{W}_k = \vect{w}_k \vect{w}_k^H \in \mathbb{H}_+^{MN_t}$，$\vect{w}_k = [\vect{w}_{1,k}^T, \ldots, \vect{w}_{M,k}^T]^T$
    \item $\mat{R}_X \triangleq \sum_{k=1}^K \mat{W}_k + \mat{Z}$
    \item $\mat{E}_m = \text{diag}(\underbrace{0,\ldots,0}_{(m-1)N_t}, \underbrace{1,\ldots,1}_{N_t}, \underbrace{0,\ldots,0}_{(M-m)N_t})$ 为 AP $m$ 选择矩阵
    \item $\mat{F}_p = \frac{2}{\sigma_s^2}\Real\{\nabla_{\boldsymbol{\theta}_p} \vect{g}_p^H \nabla_{\boldsymbol{\theta}_p}^H \vect{g}_p\} \in \mathbb{H}_+^{MN_t}$
    \item (P2-C1) 是 (P1-C1) 经 S-Procedure 处理后的 LMI 形式，详见 §逐步凸化 Step 3
\end{itemize}

\textbf{(P1) $\Leftrightarrow$ (P2) 严格等价}：变量映射 $\vect{w}_k \leftrightarrow \mat{W}_k = \vect{w}_k \vect{w}_k^H$ 为一一对应（秩一约束保证可恢复），所有约束都通过 $\tr(\vect{x}\vect{x}^H \mat{W}) = \vect{x}^H \mat{W} \vect{x}$ 等恒等式等价改写。**无任何松弛**，仅是数学形式重写。

%==============================================================================
\section{非凸性分析}
%==============================================================================

\subsection{目标函数}

目标函数 $\sum_{m,k} \|\vect{w}_{m,k}\|^2 + \sum_m \tr(\mat{Z}_m)$ 是平方范数与迹之和。每个项 $\|\vect{w}_{m,k}\|^2 = \vect{w}_{m,k}^H \vect{w}_{m,k}$ 是凸二次函数（正定 Hessian）。但与 lifted 形式 (P2-C7) 的秩一约束结合后，可行集变为非凸（rank-1 流形）。

\textbf{凸性状态}：目标本身凸，但整体问题因约束而非凸。

\subsection{通信 SINR 约束 (P1-C1)}

\begin{equation}
\frac{|\vect{h}_k^H \vect{w}_k|^2}{\sum_{j\neq k}|\vect{h}_k^H \vect{w}_j|^2 + \sigma_c^2} \geq \gamma_k
\end{equation}

这是广义分式规划约束。函数 $f(\vect{w}) = \frac{|\vect{h}^H \vect{w}|^2}{\sum |\vect{h}^H \vect{w}_j|^2 + \sigma^2}$ 单独对每个 $\vect{w}_k$ 拟凸（固定其他变量），但对 $\{\vect{w}_k\}_{k=1}^K$ 联合非凸（分子分母变量间双线性耦合）。

\textbf{非凸原因}：超水平集 $\{ \vect{w} : f(\vect{w}) \geq \gamma \}$ 非凸。对满足约束的 $\vect{w}^{(1)}, \vect{w}^{(2)}$，其凸组合 $\theta \vect{w}^{(1)} + (1-\theta) \vect{w}^{(2)}$ 通常违反 SINR 约束（分母交叉项）。

\textbf{非凸源}：NC1（SINR 分式结构）+ NC2（最坏情况 CSI 不确定性导致半无限约束）。

\subsection{感知 SINR 约束 (P1-C2)}

\begin{equation}
\frac{|\vect{g}_p^H \mat{Z} \vect{g}_p|}{\sigma_s^2} \geq \gamma_S^{\text{PoD}}
\end{equation}

对 $\mat{Z}$ 看似线性，但 $\mat{Z}$ 通过功率约束 (P1-C5) 与 $\{\vect{w}_k\}$ 耦合，使联合可行集非凸。

\subsection{PCRB 约束 (P1-C3)}

FIM 迹涉及依赖于 $\mat{R}_X$ 的矩阵逆：
\begin{equation}
\tr\left( \left[ \mat{J}_p^{\text{prior}} + \mat{J}_p^{\text{data}}(\mat{R}_X) \right]^{-1} \right) \leq \Gamma_{\text{Track},p}
\end{equation}

\textbf{非凸原因}：矩阵逆 $\mat{X}^{-1}$ 在 $\mat{X} \succ \mat{0}$ 上凸（矩阵凸），但迹的逆是凸函数。此处是凸函数的下界约束，对应凸函数的上界——这是非凸约束。

等价地，$\tr(\mat{J}_p^{\text{data}}) \geq \Gamma_{\text{Track},p}$ 对 $\mat{J}_p^{\text{data}}$ 线性，但 $\mat{J}_p^{\text{data}}$ 是 $\mat{R}_X$ 的非线性函数（信道梯度二次型）。

\textbf{非凸源}：NC6（矩阵逆 / FIM 非线性依赖）。

\subsection{AP 选择约束 (P1-C4)}

\begin{equation}
\sum_{m=1}^{M} b_{mp} = N_{\text{req}}, \quad b_{mp} \in \{0, 1\}
\end{equation}

\textbf{非凸原因}：集合 $\{0, 1\}$ 是离散两点集。凸包为 $[0, 1]$，但约束要求整数值。

\textbf{非凸源}：NC3（二进制组合）。

\subsection{功率约束 (P1-C5)}

\begin{equation}
\sum_{k=1}^{K} \|\vect{w}_{m,k}\|^2 + \tr(\mat{Z}_m) \leq P_{\max}
\end{equation}

\textbf{凸性状态}：凸约束（对 $\vect{w}_{m,k}$ 二阶锥约束，对 $\mat{Z}_m$ 线性）。但耦合通信与感知变量，与其他约束结合后贡献整体非凸性。

\textbf{状态}：已凸，无需变换。

\subsection{PSD 约束 (P1-C6)}

$\mat{Z}_m \succeq \mat{0}$ 定义凸锥（半正定锥）。

\textbf{状态}：已凸，无需变换。

\subsection{二进制约束 (P1-C7)}

$b_{mp} \in \{0, 1\}$ 与 (P1-C4) 同非凸性。

\subsection{秩一约束 (P1-C8)}

$\rank(\mat{W}_k) = 1$，其中 $\mat{W}_k = \vect{w}_k \vect{w}_k^H$。

\textbf{非凸原因}：秩一 PSD 矩阵集合为：
\begin{equation}
\mathcal{F}_{\text{rank-1}} = \{ \mat{W} \succeq \mat{0} : \rank(\mat{W}) = 1 \}
\end{equation}

这是非凸流形。例如 $\mat{W}^{(1)} = \vect{v}_1 \vect{v}_1^H$ 和 $\mat{W}^{(2)} = \vect{v}_2 \vect{v}_2^H$（$\vect{v}_1 \perp \vect{v}_2$），其平均 $\frac{1}{2}(\mat{W}^{(1)} + \mat{W}^{(2)})$ 秩为 2，违反约束。

\textbf{非凸源}：NC4（秩一要求）。

\begin{table}[htbp]
\centering
\caption{问题 (P1) 非凸性源识别}
\begin{tabular}{clll}
\toprule
ID & 源 & 影响约束 & 凸性违反 \\
\midrule
NC1 & SINR 分式结构 & (P1-C1), (P1-C2) & 二次型之比 \\
NC2 & 最坏情况 CSI 不确定性 & (P1-C1) & 半无限约束 \\
NC3 & 二进制 AP 选择 & (P1-C4), (P1-C7) & 离散集 $\{0,1\}$ \\
NC4 & 秩一要求 & (P1-C8) & 非凸流形 \\
NC5 & 变量耦合 & (P1-C1)--(P1-C5) & 双线性项 \\
NC6 & PCRB 中矩阵逆 & (P1-C3) & FIM 非线性依赖 \\
\bottomrule
\end{tabular}
\end{table}

%==============================================================================
\section{逐步凸化}
%==============================================================================

我们对 (P2) 施以六步变换，逐步将每个非凸源转化为凸形式。整体思路是：先以 SDR 牺牲 rank-1 紧致性换取可行域的凸化，再以 S-Procedure 处理半无限鲁棒约束，最后对剩余 PCRB / 感知 SINR / 功率约束做线性化或保留。每步变换都尽量保持严格等价（仅 Step 1 在 $K>2$ 时为紧致松弛，Step 6 为工程启发式），并在文中即时注明等价性。

我们首先处理非凸性中最棘手的部分——rank-1 约束 (P2-C7)。尽管 (P1)↔(P2) 严格等价保留了 rank-1，但 rank-1 流形本身是非凸的。SDR 通过将 (P2-C7) 松弛为仅 $\mat{W}_k \succeq \mat{0}$ 来去除这一非凸性：
\begin{equation}
\mat{W}_k \succeq \mat{0}, \quad \forall k. \tag{S1.1}
\end{equation}
可行域从 rank-1 流形 $\mathcal{F}_{\text{rank-1}}$ 扩大为半正定锥 $\mathcal{F}_{\text{SDR}} = \{\mat{W} \succeq \mat{0}\}$。由 Luo 等关于 MISO 多播问题的紧致性结论，$K \leq 2$ 时 SDR 紧致（最优 $\mat{W}_k^*$ 自然满足 $\rank(\mat{W}_k^*)=1$），$K > 2$ 时 SDR 提供原问题下界 $P_{\text{SDR}}^* \leq P_{\text{original}}^*$；后者情形下，高斯随机化以 $O(1/L)$ 性能损失恢复可行 rank-1 解。

处理完 rank-1 后，约束 (P2-C1) 中仍残留分式结构与最坏情况不确定性，必须按两步走：先把分式改写为 S-Procedure 可处理的二次型，再让 S-Procedure 处理 $\Delta\vect{h}$ 的不确定性。具体地，固定某个最坏情况 $\Delta\vect{h}_k$ 满足 $\|\Delta\vect{h}_k\| \leq \epsilon_h \|\hat{\vect{h}}_k\|$，SINR 不等式
$$
\frac{\tr(\hat{\mat{H}}_k \mat{W}_k) + \text{($\Delta\vect{h}$ 修正项)}}{\sum_{j\neq k} \tr(\hat{\mat{H}}_k \mat{W}_j) + \text{($\Delta\vect{h}$ 修正项)} + \sigma_c^2} \geq \gamma_k
$$
在分母恒为正（$\sigma_c^2 > 0$）的前提下，交叉相乘化为二次型不等式
\begin{equation}
\tr(\hat{\mat{H}}_k \mat{W}_k) - \gamma_k \sum_{j\neq k} \tr(\hat{\mat{H}}_k \mat{W}_j) \geq \gamma_k \sigma_c^2. \tag{S2.1}
\end{equation}
这一步是严格等价而非近似——分母为正保证了乘以分母的双向成立。得到的二次型 (S2.1) 正是 S-Procedure 的输入。

现在将最坏情况 $\Delta\vect{h}_k$ 恢复为变量。定义 $\mat{A}_k \triangleq \frac{1}{\gamma_k} \mat{W}_k - \sum_{j\neq k} \mat{W}_j$ 后，(P2-C1) 可改写为关于 $\Delta\vect{h}_k$ 的二次不等式
$$
\tilde{f}_1(\Delta\vect{h}) = (\hat{\vect{h}} + \Delta\vect{h})^H \mat{A}_k (\hat{\vect{h}} + \Delta\vect{h}) - \sigma_c^2 \geq 0, \quad \forall \|\Delta\vect{h}_k\| \leq \epsilon_h \|\hat{\vect{h}}_k\|.
$$
不确定集由二次约束 $\tilde{f}_2(\Delta\vect{h}) = \epsilon_h^2 \|\hat{\vect{h}}_k\|^2 - \|\Delta\vect{h}_k\|^2 \geq 0$ 描述。S-Procedure 指出，对范数球上的单个二次约束，``$\forall \Delta\vect{h} \in \mathcal{B}_\epsilon: \tilde{f}_1 \geq 0$'' 等价于 ``$\exists \mu_k \geq 0: \tilde{f}_1 - \mu_k \tilde{f}_2 \geq 0, \forall \Delta\vect{h}$''——后者是关于 $\Delta\vect{h}$ 的二次型 $\succeq 0$ 条件，可写为 LMI：
\begin{equation}
\exists \mu_k \geq 0: \begin{bmatrix} \mat{A}_k + \mu_k \mat{I} & \mat{A}_k \hat{\vect{h}}_k \\[1.2ex] \hat{\vect{h}}_k^H \mat{A}_k & \hat{\vect{h}}_k^H \mat{A}_k \hat{\vect{h}}_k - \sigma_c^2 - \mu_k \epsilon_h^2 \|\hat{\vect{h}}_k\|^2 \end{bmatrix} \succeq \mat{0}. \tag{S3.1}
\end{equation}
充分性显然（$\tilde{f}_1 - \mu_k \tilde{f}_2 \succeq 0$ 直接给出 $\tilde{f}_1 \geq \mu_k \tilde{f}_2 \geq 0$ 在球内）；必要性由 $\mat{A}_k + \mu_k \mat{I} \succeq \mat{0}$ 与 Lagrangian 对偶保证，参见 Luo 2010 推论 3.4。新增的 $\mu_k \geq 0$ 是 S-Procedure 松弛变量。

接下来处理感知侧与跟踪精度约束。PCRB 约束 (P2-C3) 涉及 Fisher 信息矩阵的迹——但 FIM 本身对 $\mat{R}_X$ 是\textbf{仿射}的，原因是 $\nabla_{\boldsymbol{\theta}_p} \vect{g}_p$ 在当前时隙由目标预测状态确定，可视为已知常数。在 $\nabla_{\boldsymbol{\theta}_p} \vect{g}_p \in \mathbb{C}^{MN_t \times D}$（$D$ 为目标状态维度）下，$\mat{J}_p^{\text{data}} \in \mathbb{C}^{D\times D}$，其迹通过循环性质 $\tr(\mat{A}\mat{B}\mat{C}) = \tr(\mat{C}\mat{A}\mat{B})$ 化为对 $\mat{R}_X$ 的线性函数：
\begin{align}
\tr(\mat{J}_p^{\text{data}}) &= \frac{2}{\sigma_s^2} \Real\Big\{ \tr\big(\nabla_{\boldsymbol{\theta}_p}^H \vect{g}_p \cdot \mat{R}_X \cdot \nabla_{\boldsymbol{\theta}_p} \vect{g}_p^H\big) \Big\} \notag \\
&= \frac{2}{\sigma_s^2} \Real\Big\{ \tr\big(\mat{R}_X \cdot \underbrace{\nabla_{\boldsymbol{\theta}_p} \vect{g}_p^H \cdot \nabla_{\boldsymbol{\theta}_p}^H \vect{g}_p}_{\in\, \mathbb{C}^{MN_t \times MN_t}}\big) \Big\}. \tag{S4.1a}
\end{align}
定义 Hermitian PSD 常数矩阵
$$
\mat{F}_p = \frac{2}{\sigma_s^2} \Real\{\nabla_{\boldsymbol{\theta}_p} \vect{g}_p^H \cdot \nabla_{\boldsymbol{\theta}_p}^H \vect{g}_p\} \in \mathbb{H}_+^{MN_t}
$$
（$\nabla_{\boldsymbol{\theta}_p} \vect{g}_p^H \in \mathbb{C}^{MN_t \times D}$、$\nabla_{\boldsymbol{\theta}_p}^H \vect{g}_p \in \mathbb{C}^{D \times MN_t}$，外积为 $\mathbb{C}^{MN_t \times MN_t}$；PSD 由 $\vect{x}^H(\mat{A}^H\mat{A})\vect{x} = \|\mat{A}\vect{x}\|^2 \geq 0$ 保证），则 (P2-C3) 等价于
\begin{equation}
\tr(\mat{F}_p \mat{R}_X) \geq \Gamma_{\text{Track},p}, \tag{S4.1b}
\end{equation}
即 (P2-C3) 严格等价于 (S4.1b)，后者对 $\mat{W}_k, \mat{Z}$ 是线性约束。

感知 SINR 约束 (P2-C2) 在提升域下同样简单：$\frac{|\vect{g}_p^H \mat{Z} \vect{g}_p|}{\sigma_s^2} \geq \gamma_S^{\text{PoD}}$ 直接交叉相乘（$\sigma_s^2$ 为常数，$\mat{Z} \succeq \mat{0} \Rightarrow \vect{g}_p^H \mat{Z} \vect{g}_p \geq 0$ 保证分子非负）得到
\begin{equation}
\tr(\vect{g}_p \vect{g}_p^H \mat{Z}) \geq \gamma_S^{\text{PoD}} \sigma_s^2, \tag{S4.2}
\end{equation}
对 $\mat{Z}$ 线性。

per-AP 功率约束 (P2-C4) 在提升域下保持线性：利用 $\mat{E}_m$ 选择 AP $m$ 的天线分量（$\mat{E}_m \in \mathbb{R}^{MN_t \times MN_t}$ 对角选择矩阵），$\sum_k \|\vect{w}_{m,k}\|^2 + \tr(\mat{Z}_m) = \tr(\mat{E}_m \mat{R}_X)$，故 (P2-C4) 等价于
\begin{equation}
\tr(\mat{E}_m \mat{R}_X) \leq P_{\max}, \tag{S5.1}
\end{equation}
无需任何变换。

最后处理 AP 选择约束 (P2-C8) 的离散性。该约束是 NP-hard（$K$-medoid 变种），我们采用两步分解近似：外层按大尺度衰落排序 $\text{PL}(d_{m,p})$ 选取 top-$N_{\text{req}}$ AP，确定服务集 $\mathcal{M}_p$；内层在固定 $\mathcal{M}_p$ 上求解凸 SDP。这一分解\textbf{无理论最优性保证}——外层选择可能不是真正的最优组合（NP-hard），但工程实践证明足够，且内层 SDP 在固定 AP 集上仍是凸紧的。

%==============================================================================
\section{最终凸 SDP 问题 (P3) — (P2) 的凸松弛}
%==============================================================================

综合 §逐步凸化 7 步，从 \textbf{(P2)} 出发，丢弃 (P2-C7) 秩一约束并对 (P2-C8) 启发式处理后，得到 \textbf{(P3)} 凸 SDP：

\begin{align}
\min_{\{\mat{W}_k\}, \mat{Z}, \{\mu_k\}} \quad 
& \sum_{k=1}^K \tr(\mat{W}_k) + \tr(\mat{Z}) \tag{P3} \\
\text{s.t.} \quad 
& \begin{bmatrix} \mat{A}_k + \mu_k \mat{I} & \mat{A}_k \hat{\vect{h}}_k \\[1.5ex] \hat{\vect{h}}_k^H \mat{A}_k & \hat{\vect{h}}_k^H \mat{A}_k \hat{\vect{h}}_k - \sigma_c^2 - \mu_k \epsilon_h^2 \end{bmatrix} \succeq \mat{0}, \quad \forall k \tag{P3-C1} \\
& \tr(\vect{g}_p \vect{g}_p^H \mat{Z}) \geq \gamma_S^{\text{PoD}} \sigma_s^2, \quad \forall p \tag{P3-C2} \\
& \tr(\mat{F}_p \mat{R}_X) \geq \Gamma_{\text{Track},p}, \quad \forall p \tag{P3-C3} \\
& \tr(\mat{E}_m \mat{R}_X) \leq P_{\max}, \quad \forall m \tag{P3-C4} \\
& \mat{W}_k \succeq \mat{0}, \quad \forall k \tag{P3-C5} \\
& \mat{Z} \succeq \mat{0} \tag{P3-C6} \\
& \mu_k \geq 0, \quad \forall k \tag{P3-C7}
\end{align}

其中 $\mat{A}_k = \frac{1}{\gamma_k} \mat{W}_k - \sum_{j \neq k} \mat{W}_j$，$\mat{F}_p$ 在 Step 5a 中定义，$\mat{E}_m$ 为 AP 选择矩阵。

\begin{table}[htbp]
\centering
\caption{(P1) 与 (P3) 约束对应关系}
\begin{tabular}{clll}
\toprule
(P1) 约束 & (P3) 约束 & 变换 & 凸性 \\
\midrule
(P1-C1) SINR & (P3-C1) LMI & Steps 1,2,3,4 & LMI（凸） \\
(P1-C2) 感知 SINR & (P3-C2) 线性 & Steps 1,2,5b & 线性（凸） \\
(P1-C3) PCRB & (P3-C3) 线性 & Steps 1,2,5a & 线性（凸） \\
(P1-C4) AP 选择 & (P3-C4) 固定集 & Step 7 & 线性（凸） \\
(P1-C5) 功率 & (P3-C4) 线性 & Steps 1,6 & 线性（凸） \\
(P1-C6) PSD & (P3-C6) PSD & 不变 & PSD 锥（凸） \\
(P1-C7) 二进制 & 消除 & Step 7 & 启发式 \\
(P1-C8) 秩一 & (P3-C5) PSD & Step 2 & PSD 锥（凸） \\
\bottomrule
\end{tabular}
\end{table}

\begin{theorem}[问题凸性与等价性]
问题 (P3) 是标准凸 SDP。所有约束为 LMI 或线性不等式，目标线性。$K \leq 2$ 时 (P3) $\equiv$ (P2) $\equiv$ (P1)（SDR 紧致 + (P1)↔(P2) 严格等价 + Step 1-6 等价 + Step 7 启发式不影响此结论）。$K > 2$ 时 (P3) 提供 (P2) 的下界，即 (P1) 的下界。(P3) 可通过内点法在多项式时间求解，复杂度上界为 $O\big((K+1)^3 (MN_t)^6 \cdot (K+P+M)\big)$——精确源于 (P3) 决策变量维数 $n = O(K (MN_t)^2)$，LMI 约束最大尺寸 $(MN_t+1)$；简化量级 $O((K M N_t^2)^{3.5})$。
\end{theorem}

\begin{proof}
(P3) 中各约束为 LMI (P3-C1)、线性不等式 (P3-C2, P3-C3, P3-C4) 或 PSD 锥约束 (P3-C5, P3-C6)。目标线性。故 (P3) 为凸 SDP。变换链：
\begin{itemize}
    \item \textbf{(P1) $\Leftrightarrow$ (P2)}: 严格等价（变量提升，rank-1 约束保证可恢复）
    \item \textbf{Step 1-6}：严格等价（协方差提升、SDR 紧致条件、S-Procedure 充要、分式线性化、PCRB 仿射、感知 SINR 线性化、功率约束）
    \item \textbf{Step 2 (SDR) 在 $K>2$ 时为紧致松弛}：$P_{\text{P3}}^* \leq P_{\text{P2}}^* = P_{\text{P1}}^*$
    \item \textbf{Step 7 (AP 启发式)}：外层 AP 集合由大尺度衰落排序选取，\textbf{无理论最优性保证}——既可能给出原问题下界，也可能因启发式选错 AP 而给出不可比较的目标值
\end{itemize}
故 $K \leq 2$ 时 (P3) $\equiv$ (P1)；$K > 2$ 时 $P_{\text{P3}}^* \leq P_{\text{P1}}^*$。复杂度上界由内点法对 SDP 的标准结论给出（Nesterov-Todd 路径跟踪，最坏情况 $O(n^3 L)$，其中 $n$ 为决策变量维数，$L$ 为输入数据位长）。
\end{proof}

\end{document}
